\section{Knowledge-extraction pipeline output and provenance}\label{section:results-ke-pipeline-output-and-provenance}

This section reports the outputs of the knowledge-extraction pipeline, per stage, following the evaluation setup of Subsection~\ref{subsec:eval-knowledge-extraction}. Then it measures how consistently provenance is preserved throughout the pipeline. The results in this section are computed on the frozen five-voter cascade output over the 15 related papers, since a greater number of supporting and contradicting rules was expected on a set of related papers, rather than a randomly selected subset of papers. The best configuration found during the calibration stage ($\theta = 0.9$, $\theta_{\mathrm{reject}} = 0.2$, and grounding threshold~$0.5$) was used, as explained in Section~\ref{section:results-abc-calibration}. All values were reproduced offline from the primer via a cache bridge (see the footnote), without any further API calls.\footnote{
The cost of calls to voter LLM APIs is paid only once; their outputs are gathered into a per-chunk cache (the \emph{primer}). The cache-bridge reads the voter outputs at the configuration frozen after calibration directly from this cache and reruns the deterministic downstream stages without further API calls. Configuration changes, such as dropping voters for experiments, are made using this cache.}

Two observations can be made by looking at the numbers in Table~\ref{tab:knowledge-extraction-funnel}. First, rule consolidation is limited. The NORMALIZE, GROUP, and CANONICALIZE stages retain most of the rules: 1\,747 canonical rules are produced from 1\,860 normalized findings, corresponding to approximately 93.9\%. This can be explained by the MAP stage emitting findings that are mostly atomic and distinct, correctly following the prompt instructions. Therefore, the pipeline produces approximately one final rule per normalized finding, rather than compressing many findings into a smaller set of rules. The RESOLVE stage does not change the number of canonical rules that were produced; it only scores each canonical rule and orders them based on this score.

\begin{table}[htpb]
    \caption[Knowledge-extraction funnel]{Knowledge-extraction funnel for the \texttt{related15} set: 15 papers processed with the frozen cascade.}
    \label{tab:knowledge-extraction-funnel}
    \centering
    \small
    \begin{tabular}{@{}l r l@{}}
        \toprule
        Stage & Count & Retained vs previous \\
        \midrule
        MAP findings
            & 2\,294
            & -- \\
        $\rightarrow$ grounded
            & 1\,923
            & 83.8\% \\
        $\rightarrow$ normalized/deduplicated
            & 1\,860
            & 96.7\% \\
        $\rightarrow$ groupable
            & 1\,820
            & 97.8\% \\
        $\rightarrow$ canonical rules
            & 1\,747
            & 96.0\% \\
        $\rightarrow$ final rules
            & 1\,747
            & 116.5 per paper \\
        Relations
            & 25
            & 13 intra-paper + 12 cross-paper, including 3 contradictions \\
        \bottomrule
    \end{tabular}
\end{table}

Second, the relation output is sparse. Across 15 related papers, the RELATE stage emits 13 intra-paper and 12 cross-paper relations. Only three of these 25 relations are labeled CONTRADICT: one intra-paper and two cross-paper contradictions. This sparsity explains why a synthetic balanced set of claim-pairs was used for Section~\ref{section:results-nli-in-grounding-and-relation-classification}, in order to evaluate the relation classifier: the corpus simply does not contain enough relations or contradictions to estimate per-class accuracy. The corresponding funnel on the held-out evaluation set is reported in Section~\ref{section:results-generalization-to-heldout-papers}, with the generalization strict-F1 score.

The provenance carry-rate measures whether every extracted rule or artifact can be traced back to its source. Measured on the entire ingested corpus with 977 documents, 35\,896 text elements, 4\,479 figures, and 1\,960 tables, the provenance chain is complete, and every populated field is traceable to its source. All paragraphs have a \texttt{unique\_path} field containing the PMCID, section path, and position of the paragraph within the section, making it possible to trace them back to their source papers. Provenance is complete at the rule level as well: all 1\,747 final rules scored on the \texttt{related15} cluster resolve back to at least one source paragraph from their own paper, yielding a rule-level carry-rate of 100\%.

Media provenance is also strong for table and figure objects. They can not only be traced back to the papers but also to the coordinates on the pages where they appear: both extracted figures and tables carry their cropped images and bounding-box information 100\% of the time. All extracted tables include their captions; however, the percentage of figures with captions is 86.3\%. Cross-reference linking, where figures and tables are linked to the texts in which they are mentioned, resolves for 53.6\% of figures and 58.9\% of tables. Only 12.8\% of paragraphs cite any figures or tables, but more importantly, no downstream stage in the pipeline uses the cross-reference linking outputs, so this does not affect any results in this thesis. These results are reported only for completeness. 

RQ2 also investigates the quality of the extracted findings, but the strict-F1 quality is determined by the MAP-stage cascade configuration, and is therefore reported together with the cascade calibration and cost-quality comparison (§\ref{section:results-abc-calibration}--\ref{section:results-cost-quality-comparison}): 0.7160 on \texttt{related15}. The held-out strict-F1 of $0.7128$ on \texttt{heldout15} belongs to the generalization test (RQ5, Section~\ref{section:results-generalization-to-heldout-papers}). This section, therefore, covers only the pipeline funnel and the provenance carry rate.